% Phương pháp triển khai và kịch bản. Phụ trách: TV2.
\phutrach{TV2 -- phương pháp triển khai và kịch bản}

\subsection{Phương pháp triển khai thực nghiệm}
\ghichu{Lệnh kiểm tra và \code{team\_hash}; sinh dữ liệu seed 20261116; chuyển hệ, nạp, checksum; kiểm tra bản sao; chờ đĩa lắng; chạy tiếp sau khi bị ngắt; chỉ so hai hệ trên cùng máy, giữa các máy chỉ so tỉ lệ. Cấu hình máy và hồ sơ XS ở Chương~\ref{sec:cai-dat}.}

\begin{table}[!htp]
\centering\bangfont
\caption{Đối chiếu cấu hình và dữ liệu giữa ba máy}\label{tab:checksum}
\begin{tabular}{|C{1.4cm}|L{3.6cm}|L{4.6cm}|L{4.6cm}|}
\hline
\dongdau \textbf{Máy} & \textbf{team\_hash} & \textbf{sum\_cod, 1 triệu dòng} & \textbf{sum\_cod, 3 triệu dòng} \\ \hline
M1 & \chuaco & \chuaco & \chuaco \\ \hline
M2 & fa4c12cdf33a & 5000968738439 & --- \\ \hline
M3 & \chuaco & \chuaco & \chuaco \\ \hline
\end{tabular}
\end{table}

\subsection{Kịch bản thử nghiệm}
Bảng~\ref{tab:kich-ban} tóm tắt các kịch bản. Thiết kế chi tiết nằm ở đầu mỗi chương kết quả.

\begin{table}[!htp]
\centering\bangnho
\caption{Tổng hợp các kịch bản thử nghiệm}\label{tab:kich-ban}
\begin{tabular}{|C{1.3cm}|L{3.4cm}|C{1.5cm}|L{2.4cm}|L{3.0cm}|L{2.4cm}|}
\hline
\dongdau \textbf{Mã} & \textbf{Biến thay đổi} & \textbf{Máy} & \textbf{Số dòng} & \textbf{Chỉ số} & \textbf{Trình bày} \\ \hline
B.1 & Số chỉ mục (0, 5) & M1 & 1 triệu, 3 triệu & Dung lượng, tài nguyên tạo chỉ mục & Chương~\ref{sec:dung-luong} \\ \hline
RQ1 & Covering, R, truy vấn & M3, M1 & 1 triệu (M3), 3 triệu (M1) & Kế hoạch, độ trễ, cop task, RPC & Chương~\ref{sec:doc} \\ \hline
RQ2 & Độ chọn lọc, chế độ ép kế hoạch & M3 & 3 triệu & Độ trễ, điểm giao nhau & Chương~\ref{sec:doc} \\ \hline
RQ4 & Số chỉ mục, tải, cấu hình bản sao & M2 & 1 triệu & TPS, p95, p99, byte ghi & Chương~\ref{sec:ghi} \\ \hline
\end{tabular}
\end{table}

\subsection{Cách xử lý số liệu}
\ghichu{Trung vị và min--max của các vòng; không bỏ lượt; lượt bất thường được nêu riêng; mỗi CSV ghi máy, hồ sơ, cấu hình bản sao và mã cấu hình.}
